Three ablations (three seeds each, seeds 0--2) probe the reflow and the diffusion baseline; they were chosen before running, and a fourth variant, the cosine-start diagnostic, was added after seeing the first cosine result.

\textbf{Diffusion schedule.} Replacing the linear $\beta$ schedule by the cosine schedule of~\cite{nichol2021improved} makes DDIM fail without any $\hat x_0$ clipping (Table~\ref{tab:abls}, top): even at $K{=}100$ the SW is above 1 on every task, and at $K{=}1$ it is several hundred. A diagnostic variant changes only the start index of the uniform grid (it starts at the last index with $\bar\alpha\ge4\times10^{-5}$, about the linear schedule's terminal value; same trained network): with that start index the 100-step SW drops to 0.054, 0.048 and 0.074. The failure is therefore tied to the first sampling step; it may arise because at $\bar\alpha_T\approx0$ the estimate $\hat x_0$ divides the noise-prediction error by $\sqrt{\bar\alpha}$, which we did not test separately. The truncated-start variant is nominally better than linear-schedule DDIM at $K{=}1$ (within one standard deviation on the checkerboard, and we did not test matched seeds) but worse at $K{=}16$ on all three tasks and at $K{=}4$ on eight Gaussians and two moons (Tables~\ref{tab:sw} and~\ref{tab:abls}, Fig.~\ref{fig:abl}; seeds 0--2 against 0--4 for the linear row), so, with three seeds, the weak few-step DDIM baseline in Table~\ref{tab:sw} does not appear to be an artefact of choosing the linear schedule, although we did not tune other samplers or schedules.

\begin{table*}[t]\centering\caption{Ablations (mean $\pm$ std over seeds 0--2). Top: cosine-schedule DDIM with and without the truncated start. Bottom: reflow with a 4-step teacher and a second reflow round (compare with Reflow-1 in Table~\ref{tab:sw}).}\label{tab:abls}
\footnotesize\resizebox{0.75\textwidth}{!}{\begin{tabular}{llcccccc}
\toprule
Method & Task & SW@1 $\downarrow$ & SW@4 $\downarrow$ & SW@16 $\downarrow$ & SW@100 $\downarrow$ & MMD@100 $\downarrow$ & $n$ \\
\midrule
DDIM (cosine) & 8 Gaussians & 648.562 $\pm$ 105.784 & 47.441 $\pm$ 7.932 & 10.331 $\pm$ 1.775 & 1.174 $\pm$ 0.268 & 0.102 $\pm$ 0.033 & 3 \\
DDIM (cosine, start at $\bar\alpha\ge4\times10^{-5}$) & 8 Gaussians & 2.780 $\pm$ 0.588 & 0.461 $\pm$ 0.189 & 0.242 $\pm$ 0.067 & 0.054 $\pm$ 0.004 & 0.002 $\pm$ 0.000 & 3 \\
\midrule
DDIM (cosine) & Two moons & 586.667 $\pm$ 58.955 & 62.387 $\pm$ 5.553 & 15.683 $\pm$ 2.192 & 2.142 $\pm$ 0.370 & 0.144 $\pm$ 0.021 & 3 \\
DDIM (cosine, start at $\bar\alpha\ge4\times10^{-5}$) & Two moons & 1.952 $\pm$ 0.477 & 0.347 $\pm$ 0.072 & 0.166 $\pm$ 0.050 & 0.048 $\pm$ 0.015 & 0.002 $\pm$ 0.001 & 3 \\
\midrule
DDIM (cosine) & Checkerboard & 545.477 $\pm$ 117.538 & 64.909 $\pm$ 16.613 & 16.593 $\pm$ 4.772 & 2.268 $\pm$ 0.599 & 0.184 $\pm$ 0.033 & 3 \\
DDIM (cosine, start at $\bar\alpha\ge4\times10^{-5}$) & Checkerboard & 2.145 $\pm$ 0.786 & 0.217 $\pm$ 0.037 & 0.229 $\pm$ 0.078 & 0.074 $\pm$ 0.017 & 0.005 $\pm$ 0.002 & 3 \\
\bottomrule
\end{tabular}
}\\[3pt]
\resizebox{0.75\textwidth}{!}{\begin{tabular}{llcccccc}
\toprule
Method & Task & SW@1 $\downarrow$ & SW@4 $\downarrow$ & SW@16 $\downarrow$ & SW@100 $\downarrow$ & straight $\uparrow$ & $n$ \\
\midrule
Reflow-1 (teacher 4 steps) & 8 Gaussians & 0.099 $\pm$ 0.004 & 0.083 $\pm$ 0.004 & 0.087 $\pm$ 0.007 & 0.081 $\pm$ 0.002 & 0.996 $\pm$ 0.001 & 3 \\
Reflow-2 (Euler) & 8 Gaussians & 0.033 $\pm$ 0.007 & 0.030 $\pm$ 0.002 & 0.037 $\pm$ 0.009 & 0.029 $\pm$ 0.006 & 1.000 $\pm$ 0.000 & 3 \\
\midrule
Reflow-1 (teacher 4 steps) & Two moons & 0.269 $\pm$ 0.001 & 0.256 $\pm$ 0.005 & 0.262 $\pm$ 0.005 & 0.258 $\pm$ 0.008 & 0.995 $\pm$ 0.001 & 3 \\
Reflow-2 (Euler) & Two moons & 0.025 $\pm$ 0.003 & 0.024 $\pm$ 0.004 & 0.030 $\pm$ 0.008 & 0.026 $\pm$ 0.008 & 1.000 $\pm$ 0.000 & 3 \\
\midrule
Reflow-1 (teacher 4 steps) & Checkerboard & 0.221 $\pm$ 0.006 & 0.217 $\pm$ 0.006 & 0.220 $\pm$ 0.007 & 0.216 $\pm$ 0.006 & 0.997 $\pm$ 0.001 & 3 \\
Reflow-2 (Euler) & Checkerboard & 0.031 $\pm$ 0.002 & 0.032 $\pm$ 0.001 & 0.034 $\pm$ 0.007 & 0.030 $\pm$ 0.001 & 1.000 $\pm$ 0.000 & 3 \\
\bottomrule
\end{tabular}
}\end{table*}

\textbf{Reflow teacher and second round.} Generating the reflow pairs with a 4-step instead of a 100-step teacher caps the student near the teacher's quality at every step count shown (Table~\ref{tab:abls}, bottom: SW roughly constant in $K$, and close to FM-Euler's $K{=}4$ SW in Table~\ref{tab:sw}); the straightness is still above 0.99, so a straight student can be a bad generator when its pairs come from a poor sampler. The paired differences to Reflow-1 (last block of Table~\ref{tab:tests}) show that the 4-step teacher is worse by 0.05 to 0.24 SW, while a second reflow round changes SW at $K{=}1$ and $K{=}100$ by at most 0.009: a small improvement at $K{=}1$ ($p<0.05$ on all tasks, 3 seeds) and none at $K{=}100$ (differences of at most 0.0003, positive), for one more round of pair generation and training.


\begin{figure}[t]\centering\includegraphics[width=\columnwidth]{ablations.pdf}\caption{Ablations (sliced $W_1$, mean over seeds 0--2; line style = task). Left: diffusion schedule. Right: flow and reflow variants.}\label{fig:abl}\end{figure}
